\section{Endpoint resolvents and microscopic limits}\label{sec:limits}

Proposition~\ref{prop:edge-moments} describes finitely many eigenvectors
sampled at edges. We first turn this information into a resolvent estimate
at large fixed microscopic heights. A two-edge switch supplies the
self-consistency equation needed to descend from the height at which the
local law applies. We then construct microscopic subsequential limits and
identify both their Gaussian spectral weights and the real constants in
their resolvent representations.

\subsection{Resolvents at sampled endpoints}

Conditionally on the graph, sample an infinite sequence of independent
uniform oriented edges. For fixed \(k\), let \(W_{n,k}\) be the \(n\)-by-\(2k\)
matrix whose columns are the coordinate vectors of the first \(k\)
starting vertices, followed by those of the corresponding ending
vertices. Define
\begin{equation}\label{eq:lim-endpoint-resolvent}
 M_{n,k}(z)
 =W_{n,k}^{\mathsf T}(A-E-z/n)^{-1}W_{n,k}
 =\sum_i\frac{b_i b_i^{\mathsf T}}{x_i-z},
 \qquad b_i=W_{n,k}^{\mathsf T}u_i,\qquad \Im z>0.
\end{equation}
The normalization \(\|u_i\|=\sqrt n\) accounts for the absence of a factor
of \(n\) in the last expression. Both \(m_n\) and every real quadratic
form of \(M_{n,k}\) have nonnegative imaginary part.

For every fixed total degree, the joint moments of the \(b_i\) with
\(|x_i|\le S^2\) agree, to polynomial precision, with those of independent
centered Gaussian vectors of covariance \(C_k\). Indeed, take the distinct
eigenvector indices occurring in a monomial as the colors in
Proposition~\ref{prop:edge-moments}, and then use the conditional
independence of the sampled edges. Only finitely many colors are needed
for a fixed monomial. All errors are uniform in the selected band
indices on events of graph probability tending to one.

\begin{proposition}[Endpoint control]\label{prop:endpoint-control}
There is a deterministic \(\xi>0\), depending only on \(d,E\), with the
following property. For every fixed \(k\ge1\) and \(D>0\), there are
constants \(C_{k,D}<\infty\) and \(y_{k,D}<\infty\) such that, for every
fixed \(s\in\R\) and every fixed \(y\ge y_{k,D}\),
\begin{equation}\label{eq:lim-endpoint-control}
 \limsup_{n\to\infty}
 \PP\left(
 |m_n(s+\mathrm i y)-m_*|
 +\|M_{n,k}(s+\mathrm i y)-K_k(m_*)\|>y^{-\xi}
 \right)
 \le C_{k,D}y^{-D}.
\end{equation}
The limit is through admissible graph sizes. The constants and the height
threshold are independent of \(s\); the argument \(s+\mathrm i y\) is
fixed before taking the limit.
\end{proposition}

\begin{proof}
It suffices to prove the statement for \(k\ge2\), since the smaller
matrices are compressions. Throughout this proof \(k,s,y\) are fixed.
The exponents in polynomial errors may depend on fixed moment orders.

\paragraph{Initial estimates.}
Put \(z_0=s+\mathrm i S\). Proposition~\ref{prop:spectral-input} and
Lemma~\ref{lem:tree-identities} give, for some \(c>0\),
\begin{equation}\label{eq:lim-initial-estimates}
 \begin{split}
  m_n(z_0)&=m_*+O(S^{-c}),\\
  M_{n,k}(z_0)&=K_k(m_*)+O(S^{-c}),\\
  \Im m_n(\mathrm i S^2)+\|\Im M_{n,k}(\mathrm i S^2)\|&=O(1)
 \end{split}
\end{equation}
with probability tending to one. To justify the off-block entries in the
second assertion, write \(R=(A-E-z_0/n)^{-1}\). For independent uniform
vertices \(v,w\), the resolvent identity gives
\[
 \EE_{v,w}|R_{vw}|^2
 =\frac1{n^2}\Tr R^*R
 =\frac{\Im m_n(z_0)}{S}.
\]
Endpoints from distinct edge samples have this independent uniform
distribution. Markov's inequality makes their resolvent entries
\(O(S^{-c})\), after decreasing \(c\). Within a single sampled edge the
typical-star estimate applies, with adjacent entry
\((1+E m_*)/d\). The same reasoning gives the last assertion of
\eqref{eq:lim-initial-estimates}. These estimates also imply that the
number of indices in \(|x_i|\le S^2\) is \(O(S^2)=n^{o(1)}\).

\paragraph{Concentration under a trace bound.}
Descend from \(S\) to \(y\) by halving the height, shortening the last
step if necessary, and keep real part \(s\). Let \(\mathcal Z_n\) denote
these \(O(\log S)\) arguments. For every fixed \(C<\infty\), we claim
that, outside an event with limiting upper probability
\(O_{k,D,C}(y^{-D})\), simultaneously for \(z\in\mathcal Z_n\),
\begin{equation}\label{eq:lim-concentration-implication}
 |m_n(z)|\le C
 \quad\Longrightarrow\quad
 M_{n,k}(z)=K_k(m_n(z))+O((\Im z)^{-c_2}).
\end{equation}
Here \(c_2>0\) can be chosen independently of \(D,k,s,C\), with constants
in the error allowed to depend on the fixed parameters.

We compare \(z\) with \(z_0\). Once \(S\) exceeds a fixed multiple of
\(|s|+y+1\), for \(|x|>S^2\) we have
\begin{equation}\label{eq:lim-band-tail}
 \left|\frac1{x-z}-\frac1{x-z_0}\right|
 \le \frac{C S}{x^2}
 \le \frac C S\,\Im\frac1{x-\mathrm i S^2}.
\end{equation}
Thus the scalar and matrix contributions outside the band are
\(O(S^{-1})\) on the event in \eqref{eq:lim-initial-estimates}. For the
matrix bound, use that each residue \(b_i b_i^{\mathsf T}\) is positive
semidefinite.

Inside the band, let
\[
 w_i(z)=\frac1{x_i-z}-\frac1{x_i-z_0}.
\]
If the scalar trace at \(z\) is bounded, then
\begin{equation}\label{eq:lim-weight-squares}
 \sum_{|x_i|\le S^2}|w_i(z)|^2
 \le
 2\frac{\Im m_n(z)}{\Im z}
 +2\frac{\Im m_n(z_0)}S
 \le \frac{C'}{\Im z}.
\end{equation}
Replace the marks for this calculation by independent Gaussians of
covariance \(C_k\). For every fixed integer \(j\), each entry of
\[
 \sum_{|x_i|\le S^2}
 (b_i b_i^{\mathsf T}-C_k)w_i(z)
\]
then has \(2j\)-th absolute moment at most
\(C_{j,k,C}(\Im z)^{-j}\). In the moment expansion a spectral index
occurring only once has zero expectation. For every remaining
partition of the indices, bounded Gaussian moments and
\eqref{eq:lim-weight-squares} give the stated bound.

The same estimate, with an additional polynomially small error, holds
conditionally on each graph satisfying the required fixed-order moment
estimates. There are only \(n^{o(1)}\) terms in the expansion, and the
weights are bounded because \(\Im z\ge y\ge1\). The event
\(|m_n(z)|\le C\) depends only on the graph, so it causes no difficulty
in this conditional calculation. Choose
\[
 0<c_2<\min(c,1/4).
\]
Markov's inequality at threshold \((\Im z)^{-c_2}\) gives a failure
probability bounded by a constant times
\((\Im z)^{-j(1-2c_2)}\), plus a polynomially small error in \(n\).
Summing over the geometric sequence of heights is bounded by
\(C_{j,k,C}y^{-j/2}+o(1)\). Taking \(j\) sufficiently large proves
\eqref{eq:lim-concentration-implication}. The graph exceptional events
and the failure of \eqref{eq:lim-initial-estimates} have probability
tending to zero. Factors from \(S\) and the number of steps are
\(n^{o(1)}\), so the polynomial moment errors remain negligible.

\paragraph{A switch gives the scalar equation.}
The graph/list bijection below uses the measure-preserving switching
mechanism of \citet[Section~7.2, Propositions~7.3--7.4]{BauerschmidtHuangYauLocal}.
Figure~\ref{fig:two-edge-switch} illustrates the two-edge operation.
\begin{figure}[tbp]
 \centering
 \begin{tikzpicture}[x=1cm,y=1cm,
 vertex/.style={circle,fill=black,inner sep=2pt},
 edge/.style={line width=.8pt},every node/.style={font=\small}]
 \foreach \offset/\tag in {0/a,5.1/b}{
   \node[vertex,label=left:{$v_1$}] (\tag v1) at (\offset,1.6) {};
   \node[vertex,label=left:{$v_2$}] (\tag v2) at (\offset,0) {};
   \node[vertex,label=right:{$w_1$}] (\tag w1) at (\offset+2,1.6) {};
   \node[vertex,label=right:{$w_2$}] (\tag w2) at (\offset+2,0) {};
 }
 \draw[edge] (av1)--(aw1) (av2)--(aw2);
 \draw[edge] (bv1)--(bw2) (bv2)--(bw1);
 \draw[-{Stealth[length=2.5mm]},line width=.7pt] (2.6,.8)--(4.5,.8)
    node[midway,above] {switch};
 \node at (1,-.65) {original pairs};
 \node at (6.1,-.65) {new pairs};
\end{tikzpicture}

 \caption{A two-edge switch on an induced matching. The marked pairs
 \((v_1,w_1)\) and \((v_2,w_2)\) are replaced by
 \((v_1,w_2)\) and \((v_2,w_1)\). Each endpoint loses one edge and gains
 one edge, so its degree is preserved. All other graph edges remain in
 place, and the marked list records the new pairs. The same exchange
 reverses the operation.}
 \label{fig:two-edge-switch}
\end{figure}
With probability \(1-O_k(n^{-1})\), the \(2k\) sampled endpoints are
distinct and the only graph edges among them are the sampled pairs.
This follows directly from fixed degree and conditional independence
of the edge samples. Restrict temporarily to this event. Exchange the
ending vertices of the first two pairs and update their entries in the
edge list. The operation is a bijection of the restricted graph/list
space, and all its elements have the same probability. In particular
the switched data, denoted by hats, satisfy
\eqref{eq:lim-concentration-implication} with the same failure bound.
The restriction itself is preserved by the switch.

The adjacency update has bounded rank. The eigenvalue counting
functions before and after the update consequently differ by a bounded
amount. Integrating their difference against the derivative of
\((t-z)^{-1}\) gives
\begin{equation}\label{eq:lim-rank-trace}
 |\widehat m_n(z)-m_n(z)|\le \frac C{\Im z}.
\end{equation}
This calculation is in the microscopic spectral variable; equivalently,
\(m_n(z)=\Tr(n(A-E)-z)^{-1}\).

Write the first two oriented pairs as \((v_1,w_1),(v_2,w_2)\), and put
\[
 f_1=\frac{e_{v_1}-e_{v_2}}{\sqrt2},\qquad
 f_2=\frac{e_{w_1}-e_{w_2}}{\sqrt2},\qquad F=(f_1,f_2).
\]
The switch adds \(F B F^{\mathsf T}\), where
\[
 B=\begin{pmatrix}0&-2\\-2&0\end{pmatrix}.
\]
Let \(U=F^{\mathsf T}(A-E-z/n)^{-1}F\), and define \(\widehat U\) using
the switched graph but the same old columns \(F\). The resolvent
identity yields
\begin{equation}\label{eq:lim-switch-identity}
 (I+UB)\widehat U=U.
\end{equation}
Suppose now that \(m=m_n(z)\) lies in a fixed bounded range.
Equations~\eqref{eq:lim-concentration-implication} and
\eqref{eq:lim-rank-trace}, applied to both data sets with a sufficiently
large trace threshold, give
\begin{equation}\label{eq:lim-switch-compressions}
 U=\begin{pmatrix}m&q\\q&m\end{pmatrix}
       +O((\Im z)^{-c_2}),\qquad
 \widehat U=\begin{pmatrix}m&-q\\-q&m\end{pmatrix}
       +O((\Im z)^{-c_2}),\qquad
 q=\frac{1+E m}{d}.
\end{equation}
The sign in the second matrix comes from expressing the updated edge
list in the old columns: the right difference column is negated.
The off-diagonal entries of \eqref{eq:lim-switch-identity} now imply
\begin{equation}\label{eq:lim-scalar-equation}
 q^2-q-m^2=O((\Im z)^{-c_2}).
\end{equation}
As a quadratic in \(m\), its exact left side has simple conjugate roots
in the strict bulk. By Lemma~\ref{lem:tree-identities}, the upper root is
\(m_*\). The lower root stays a fixed distance from every \(m\) with
\(\Im m\ge0\). Hence, throughout the bounded range,
\begin{equation}\label{eq:lim-scalar-stability}
 m_n(z)=m_*+O((\Im z)^{-c_2}).
\end{equation}

\paragraph{Closing the descent.}
There is no a priori trace assumption at a lower step. Instead, if
\(\eta/2\le\eta'\le\eta\), the resolvent sum gives
\begin{equation}\label{eq:lim-halving-bound}
 \begin{split}
 |m_n(s+\mathrm i\eta')-m_n(s+\mathrm i\eta)|
 &\le \frac{\eta-\eta'}{\eta'}\Im m_n(s+\mathrm i\eta)\\
 &\le \Im m_n(s+\mathrm i\eta).
 \end{split}
\end{equation}
Indeed
\(|x-s-\mathrm i\eta'|\ge(\eta'/\eta)|x-s-\mathrm i\eta|\)
term by term. Choose a fixed trace range containing the initial value
and every possible next-step value in \eqref{eq:lim-halving-bound}
when the preceding value is within \(1\) of \(m_*\). Enlarge it once
more for the switched trace using \eqref{eq:lim-rank-trace}. Finally
take the terminal height \(y\) sufficiently large that
\eqref{eq:lim-scalar-stability} returns a value within \(1\) of \(m_*\).
Induction down the heights proves the required trace bounds and the
endpoint estimates at every step.

Combining the two concentration failure bounds, and decreasing the
final exponent from \(c_2\) to a smaller \(\xi>0\), absorbs the fixed
constants into the threshold \(y^{-\xi}\). This proves
\eqref{eq:lim-endpoint-control}. All constants are independent of \(s\):
the only additional restriction was that \(S\) eventually exceed a
fixed multiple of \(|s|+y+1\), which holds before the limit for every
fixed \(s,y\).
\end{proof}

\subsection{The limiting marked process}

We now take subsequential limits in \(n\). The probability estimates
above will be used again only after this limit has been taken. For a
locally finite point measure \(\nu=\sum_j\delta_{x_j}\), a principal
value below always means the symmetric position truncation
\(\lim_{R\to\infty}\sum_{|x_j|\le R}\).

\begin{proposition}[Microscopic limits]\label{prop:microscopic-limit}
From every sequence of admissible sizes tending to infinity one can
extract a subsequence along which
\[
 \bigl(\nu_n,m_n,(M_{n,k})_{k\ge1}\bigr)
 \ \Longrightarrow\
 \bigl(\nu,m,(M_k)_{k\ge1}\bigr)
\]
jointly, in the local vague topology for the point measures and the
compact-open topology on the upper half-plane for the analytic
functions. The limiting point measure \(\nu=\sum_j\delta_{x_j}\) is
locally finite, with multiplicities permitted. There are deterministic
\(\xi,\chi>0\), depending only on \(d,E\), such that the following
properties hold.

For every fixed \(0<\delta<1\), almost surely for all sufficiently
large \(R\),
\begin{equation}\label{eq:mesoscopic-limit}
 m(e+\mathrm i u)=m_*+O(u^{-\xi}),
 \qquad |e|\le R,\quad R^\delta\le u\le R^2,
\end{equation}
and, almost surely,
\begin{equation}\label{eq:discrepancy}
 \nu([v,w])=\rho(w-v)+O(R^{1-\chi}),
 \qquad -R\le v\le w\le R,\quad R\longrightarrow\infty.
\end{equation}
The constants and starting values of \(R\) in these almost sure
statements may depend on the sample and on \(\delta\), but the exponents
are deterministic.

On a possible extension of the probability space, for every fixed
\(k\) there are standard Gaussian vectors \(g_j\in\R^{2k}\), independent
over \(j\) conditionally on \(\nu\), such that
\begin{equation}\label{eq:marked-representation}
 \begin{split}
  m(z)&=m_{\rm r}+\PV\sum_j\frac1{x_j-z},\\
  M_k(z)&=K_k(m_{\rm r})
       +\PV\sum_j
       \frac{C_k^{1/2}g_jg_j^{\mathsf T}C_k^{1/2}}{x_j-z}.
 \end{split}
\end{equation}
These representations are consistent as \(k\) increases by retaining
the first \(k\) edge channels. A separate Gaussian mark is assigned to
each occurrence of a repeated point. The principal values converge
locally uniformly on the upper half-plane, and
\(\Im M_k(z)\) is positive definite there, almost surely simultaneously
for every finite \(k\) and every \(z\) with \(\Im z>0\).
\end{proposition}

\begin{proof}
\leavevmode
\paragraph{Tightness.}
For fixed \(L,y>0\), positivity gives
\[
 \nu_n([-L,L])\le \frac{L^2+y^2}{y}\Im m_n(\mathrm i y).
\]
Proposition~\ref{prop:endpoint-control}, with \(y\) a sufficiently
large fixed number, therefore proves tightness of the local counts.
It also proves compact-open tightness of the analytic functions.
The choice of \(y\) is made separately for each probability tolerance:
choose it so that the bound \(C_{k,D}y^{-D}\) is below that tolerance,
and then take \(n\) large. A single height is not required to handle
all tolerances. Finitely many initial sizes do not affect tightness.
Here is a direct way to see the needed compact bounds. For a scalar
resolvent sum \(f\) with nonnegative residues, and a compact set
\(\mathcal K\) in the upper half-plane,
\[
 \Im f(z)\le C_{\mathcal K,y}\Im f(\mathrm i y),
 \qquad
 |f(z)-f(\mathrm i y)|
 \le C_{\mathcal K,y}\Im f(\mathrm i y),
 \qquad z\in\mathcal K.
\]
Both inequalities follow by comparing the corresponding kernels on
the real line. Apply them to \(m_n\), and to the real quadratic forms
of \(M_{n,k}\); polarization controls the matrix entries. Thus bounded
values at \(\mathrm i y\) give uniform bounds on each fixed compact,
and Cauchy's estimates give equicontinuity on smaller compacts.

We may consequently extract the asserted joint limit by tightness,
first on growing compact sets and finitely many channel dimensions,
and then diagonally. Local vague limits with locally bounded counts
are point measures with integer multiplicities. Limits of the
analytic functions have nonnegative imaginary parts in the same
scalar or matrix sense.

\paragraph{Conditional Gaussian marks on compact intervals.}
Choose nested compact position intervals exhausting \(\R\), with
boundaries almost surely uncharged by \(\nu\). Such deterministic
boundaries can be chosen because the set of points with positive
probability of carrying an atom is countable. On each interval,
truncate its point count at a fixed integer. Order the prelimit
indices using graph data alone, resolving ties by the chosen
eigenbasis ordering. The positions and the corresponding marks are
tight: their number is bounded, and the conditional second-moment
estimates of Proposition~\ref{prop:edge-moments} control the marks.
Extract these marked lists together with the preceding limit.

For completeness, fixed-order moment matching suffices here even
though the selected indices depend on the graph. For any fixed
number of indices, the approximation is uniform over their choices.
By a diagonal selection of the graph exceptional events, one may
require that an increasing finite collection of moment orders match
with errors tending to zero, on events of probability tending to one.
More explicitly, for each fixed list length and \(k\), choose integers
\(J_n\to\infty\) sufficiently slowly that all orders at most \(J_n\)
have errors at most \(J_n^{-1}\) on a common graph event of probability
at least \(1-J_n^{-1}\). This uses only finitely many fixed-order
assertions at each stage, not a uniform growing-order estimate.
If the conditional mark laws failed to approach their product
Gaussian law in a metric for weak convergence, one could select a
deterministic sequence of graphs and index lists with this failure
while every fixed moment converged. Tightness and moment
determinacy of the Gaussian law contradict that selection.
Thus conditional weak convergence is uniform over these bounded
lists.

Bounded Lipschitz tests, first with the list length truncated, now
show that the same joint limit is obtained by assigning independent
Gaussian marks to the prelimit positions independently of the graph.
Letting the truncation increase identifies the conditional marking
law given the full unmarked process \(\nu\). The covariance in
\(2k\) endpoint coordinates is \(C_k\). At a repeated limiting
position the list can be retained on an extension of the probability
space, with a separate independent mark for each multiplicity;
its summed outer product does not require an intrinsic ordering of
the tied points. Diagonal extraction in the intervals and in \(k\)
gives consistent marks
\begin{equation}\label{eq:lim-gaussian-marks}
 b_j=C_k^{1/2}g_j
\end{equation}
for all finite channel dimensions.

\paragraph{Removing imaginary tails.}
The convergence of compact marked sums alone does not yet identify
the full resolvents. For fixed \(z\) in the upper half-plane and
large \(L\), the scalar kernels satisfy, when \(|x|>L\),
\begin{equation}\label{eq:lim-tail-domination}
 \Im\frac1{x-z}
 +\left|\frac1{x-z}-\frac1{x-\mathrm i}\right|
 \le \frac{C_z}{L}\Im\frac1{x-\mathrm i L}.
\end{equation}
The constants can be uniform for \(z\) in a fixed compact set.
For matrix sums use positivity of the residues and bound the norm
by the trace of the corresponding positive imaginary part.
Proposition~\ref{prop:endpoint-control} shows that these traces at
\(\mathrm i L\) are tight as \(L\) tends to infinity, after taking
the upper limit in \(n\). The prefactor \(L^{-1}\) in
\eqref{eq:lim-tail-domination} therefore removes the tails in
probability. Compact marked sums already converge, so
\begin{equation}\label{eq:lim-difference-identities}
 \begin{split}
  \Im m(z)&=\sum_j\Im\frac1{x_j-z},\\
  m(z)-m(\mathrm i)
    &=\sum_j\left(\frac1{x_j-z}-\frac1{x_j-\mathrm i}\right),\\
  \Im M_k(z)&=\sum_j b_jb_j^{\mathsf T}
                            \Im\frac1{x_j-z},\\
  M_k(z)-M_k(\mathrm i)
    &=\sum_j b_jb_j^{\mathsf T}
               \left(\frac1{x_j-z}-\frac1{x_j-\mathrm i}\right).
 \end{split}
\end{equation}
One may establish these identities first at countably many
arguments and then use positivity and continuity. The imaginary
sums and the sums of absolute values of the difference kernels are
finite. In particular no imaginary spectral mass, including a
linear Herglotz term, is lost at infinity. Real additive constants
are not excluded by this argument and will be identified below.

\paragraph{Uniform bounds after the microscopic limit.}
At every fixed \(s+\mathrm i y\),
Proposition~\ref{prop:endpoint-control} passes to the limit by the
portmanteau inequality for the open event where its error exceeds
the stated threshold. Its constants remain independent of the
fixed real part \(s\). Decrease \(\xi\) if necessary so that
\(\xi<1\).

Fix \(0<\delta<1\). In a padded region
\[
 |e|\le 2R,\qquad \tfrac12 R^\delta\le u\le 2R^2
\]
take a unit grid. There are \(O(R^3)\) grid points. For each fixed
dyadic \(R\), this is a finite collection of arguments in the exact
limiting law. The failure probability from
\eqref{eq:lim-endpoint-control} is at most
\(C R^{3-\delta D}\). Choose \(D\) sufficiently large and apply
Borel--Cantelli over dyadic \(R\).
The kernel identities in \eqref{eq:lim-difference-identities} give
\[
 |m'(z)|\le \frac{\Im m(z)}{\Im z}.
\]
Imaginary parts at arguments within bounded distance and comparable
height are comparable by their positive kernels. The error in
filling the unit grid is consequently \(O(1/u)\). Padding allows
all sufficiently large \(R\), not just dyadic values. This proves
\eqref{eq:mesoscopic-limit}. The limit in \(n\) preceded the growing
grid argument; no finite-\(n\) uniform estimate on expanding real
windows has been used. A countable set of \(\delta\)'s implies all
the stated fixed-\(\delta\) conclusions, by slightly decreasing
\(\delta\) when necessary.

\paragraph{From the Poisson transform to counts.}
Equation~\eqref{eq:mesoscopic-limit} first gives
\(\nu([-R,R])=O(R)\), by evaluating \(\Im m(\mathrm i R)\).
For \(-R\le v\le w\le R\), integrate
\(\pi^{-1}\Im m(e+\mathrm i\sqrt R)\) over \(e\in[v,w]\).
The result is
\[
 \rho(w-v)+O(R^{1-\xi/2}).
\]
To compare it with the sharp count, put
\[
 P_{v,w}(x)=\frac1\pi
   \int_v^w\frac{\sqrt R}{(x-e)^2+R}\dd e.
\]
The number of points within distance \(R^{3/4}\) of either
endpoint is \(O(R^{3/4})\): use
\eqref{eq:mesoscopic-limit} at that endpoint and height
\(R^{3/4}\). For points of \([-2R,2R]\) outside these two strips,
\[
 |P_{v,w}(x)-\boldsymbol 1_{[v,w]}(x)|=O(R^{-1/4}).
\]
Their total error is \(O(R^{3/4})\), since the number of points is
\(O(R)\). Finally, for \(|x|>2R\),
\(P_{v,w}(x)\le C R^{3/2}/x^2\); the linear count bound and dyadic
summation make this contribution \(O(\sqrt R)\).
These estimates are uniform in \(v,w\), and prove
\eqref{eq:discrepancy} for any sufficiently small deterministic
\[
 0<\chi<\min(1/4,\xi/2).
\]
The estimates also show that endpoint conventions have no effect
on the stated error bound.

At this point the limiting positions have a deterministic density
with a power-saving count error, and all finite spectral residues
have been identified. It remains to determine the real constants
that the compact convergence did not see.

\paragraph{Scalar principal values and their constant.}
Partial summation in \eqref{eq:discrepancy} gives, almost surely,
\begin{equation}\label{eq:lim-position-tails}
 \PV\sum_{|x_j|>R}\frac1{x_j}=O(R^{-\chi}),
 \qquad
 \sum_{|x_j|>R}\frac1{x_j^2}=O(R^{-1}).
\end{equation}
Only large \(R\) is involved, so possible points at zero do not
enter these expressions. To see the first assertion, apply the
count discrepancy separately on the positive and negative
half-lines; their leading density terms cancel in the symmetric
sum, and integration of \(O(t^{1-\chi})/t^2\) leaves
\(O(R^{-\chi})\). The second assertion follows directly from the
linear count bound.

Thus
\[
 S_\nu(z)=\PV\sum_j\frac1{x_j-z}
\]
exists locally uniformly and is analytic. Comparison with the
constant-density measure gives
\begin{equation}\label{eq:lim-pv-large-height}
 S_\nu(\mathrm i R)=\mathrm i\pi\rho+o(1).
\end{equation}
Indeed, on each half-line the counting error is
\(O(1+|v|^{1-\chi})\). Partial summation bounds its contribution
by a constant times
\[
 \int_0^\infty\frac{1+v^{1-\chi}}{v^2+R^2}\dd v=o(1),
\]
whereas the symmetric transform of constant density is
\(\mathrm i\pi\rho\).
By \eqref{eq:lim-difference-identities}, \(m-S_\nu\) is a constant
function, and its imaginary part is zero. Equations
\eqref{eq:mesoscopic-limit} and \eqref{eq:lim-pv-large-height}
identify that real constant as \(m_{\rm r}\).

\paragraph{Marked principal values and their constant.}
Fix \(k\) and condition on \(\nu\). Set
\[
 Z_j=b_jb_j^{\mathsf T}-C_k.
\]
These matrices are independent and centered, with bounded
second moments depending only on \(k,d,E\). The series
\(\sum_j Z_j/(x_j-\mathrm i)\), ordered by increasing \(|x_j|\),
converges almost surely entrywise by square summability and
independence. Also, conditionally almost surely,
\[
 \sum_{|x_j|>1}\frac{\|Z_j\|}{x_j^2}<\infty,
\]
because its conditional expectation is finite by
\eqref{eq:lim-position-tails}. Consequently the corresponding
difference series between \(z\) and \(\mathrm i\) converges
absolutely and locally uniformly. Together with the scalar
principal value this proves local uniform convergence of
\[
 S_{k,\nu}(z)=\PV\sum_j\frac{b_jb_j^{\mathsf T}}{x_j-z}.
\]

The linear count bound also gives
\[
 \sum_j|x_j-\mathrm i R|^{-2}=O(R^{-1}).
\]
Thus every entry of the centered marked series at
\(\mathrm i R\) has conditional second moment \(O(R^{-1})\),
with a sample-dependent constant. It follows that
\[
 S_{k,\nu}(\mathrm i R)
 \longrightarrow \mathrm i\pi\rho\,C_k
 \quad\hbox{in probability}.
\]
The limiting version of Proposition~\ref{prop:endpoint-control}
gives \(M_k(\mathrm i R)\to K_k(m_*)\) in probability. On the
other hand \eqref{eq:lim-difference-identities} shows that
\(M_k-S_{k,\nu}\) is almost surely constant in \(z\).
That constant must therefore equal
\[
 K_k(m_*)-\mathrm i\pi\rho\,C_k=K_k(m_{\rm r})
\]
almost surely. This proves both identities in
\eqref{eq:marked-representation}.

Finally, \eqref{eq:discrepancy} supplies infinitely many points,
and \(C_k\) is positive definite. Any \(2k\) independently
Gaussian marks span \(\R^{2k}\) almost surely. Their positive
contributions to the imaginary-part sum in
\eqref{eq:lim-difference-identities} make \(\Im M_k(z)\) positive
definite for every \(z\) in the upper half-plane on the same
event. Countability of the channel dimensions completes the
simultaneous assertion.
\end{proof}
